\chapter{Cơ sở lý thuyết}

\section{Tổng quan về ứng dụng học máy trong việc dự đoán dữ liệu chăm sóc sức khỏe}

Trong những năm gần đây, việc ứng dụng học máy trong lĩnh vực chăm sóc sức khỏe đã trở thành một xu hướng nổi bật, mang lại nhiều lợi ích trong việc dự đoán và chẩn đoán các bệnh lý. Học máy, với khả năng xử lý và phân tích dữ liệu lớn, đã được sử dụng để xây dựng các mô hình dự đoán chính xác về tình trạng sức khỏe của bệnh nhân, từ đó hỗ trợ các quyết định lâm sàng và cải thiện chất lượng chăm sóc.

\section{Đặc tính của dữ liệu y tế}
Dữ liệu y tế không chỉ lớn về quy mô mà còn phức tạp về cấu trúc, đặt ra nhiều thách thức đặc thù cho các mô hình học máy. Trong số đó, ba đặc tính dưới đây được xác định là các vấn đề trọng tâm cần giải quyết trong phạm vi khóa luận này, xuất phát trực tiếp từ đặc trưng của hai bộ dữ liệu được sử dụng: số chiều cao và tính phi tuyến là vấn đề chủ đạo của bộ Breast Cancer Gene Expression, trong khi mất cân bằng lớp và tính phi tuyến là đặc trưng nổi bật của bộ Diabetes Health Indicators. Mỗi đặc tính được trình bày theo cấu trúc: định nghĩa, hậu quả đối với mô hình và hướng xử lý trong đề tài.

\subsection{Số chiều cao}

Dữ liệu số chiều cao là dữ liệu có số lượng đặc trưng lớn hơn đáng kể so với số lượng mẫu. Trong lĩnh vực y tế, đây là đặc điểm phổ biến của nhiều loại dữ liệu sinh học, điển hình là dữ liệu biểu hiện gen trong đó mỗi mẫu bệnh nhân được mô tả bởi hàng chục nghìn giá trị biểu hiện gen trong khi số lượng bệnh nhân trong nghiên cứu thường chỉ ở mức vài trăm. Tình trạng này dẫn đến hiện tượng "lời nguyền đa chiều" (curse of dimensionality),  khi số chiều tăng, không gian đặc trưng trở nên cực kỳ thưa thớt, khoảng cách giữa các điểm dữ liệu dần mất ý nghĩa thống kê, khiến nhiều thuật toán dựa trên khoảng cách như KNN suy giảm hiệu năng rõ rệt. Đồng thời, khi số đặc trưng vượt xa số mẫu, mô hình có quá nhiều bậc tự do so với lượng thông tin có sẵn, dẫn đến nguy cơ overfitting cao, tức mô hình học thuộc dữ liệu huấn luyện thay vì nắm bắt quy luật thực sự, làm giảm khả năng tổng quát hóa trên dữ liệu mới. Ngoài ra, nhiều đặc trưng trong không gian số chiều cao thường là nhiễu hoặc dư thừa, không mang thông tin phân biệt có ý nghĩa, gây khó khăn cho quá trình học của mô hình.\\

Trong phạm vi khóa luận, nhóm áp dụng các phương pháp lựa chọn đặc trưng để giữ lại các đặc trưng thực sự có giá trị phân biệt, loại bỏ nhiễu và đặc trưng dư thừa. Điểm nổi bật của phương pháp lựa chọn đặc trưng là nó giữ nguyên các đặc trưng gốc mà không biến đổi chúng, qua đó bảo toàn khả năng giải thích của từng đặc trưng, đặc biệt quan trọng trong bối cảnh y tế khi cần truy xuất ý nghĩa sinh học của từng gene được chọn. Chi tiết các phương pháp được trình bày tại mục 2.4.

\subsection{Mất cân bằng lớp}
\label{sec:matcanbang}
Mất cân bằng lớp là hiện tượng số lượng mẫu giữa các lớp chênh lệch đáng kể, trong đó lớp thiểu số chiếm tỉ lệ rất nhỏ so với lớp đa số. Trong dữ liệu y tế, đây là đặc điểm phổ biến xuất phát từ sự chênh lệch tự nhiên giữa tỉ lệ mắc bệnh và tỉ lệ không mắc bệnh trong dân số, lớp dương tính (bệnh nhân mắc bệnh) thường chiếm tỉ lệ nhỏ hơn nhiều so với lớp âm tính (người bình thường). Khi dữ liệu mất cân bằng, các mô hình học máy có xu hướng thiên vị về phía lớp đa số vì dự đoán toàn bộ mẫu là lớp đa số đã cho độ chính xác (Accuracy) tổng thể cao. Hệ quả là mô hình thất bại trong việc nhận diện đúng các trường hợp thuộc lớp thiểu số, vốn là mục tiêu quan trọng nhất trong bối cảnh y tế: bỏ sót một ca bệnh thường gây hậu quả nghiêm trọng hơn nhiều so với cảnh báo nhầm. Điều này khiến Accuracy trở thành độ đo không phù hợp khi dữ liệu mất cân bằng, và các độ đo như F1-score và ROC-AUC cần được ưu tiên thay thế.\\

Trong phạm vi khóa luận, nhóm áp dụng hai kỹ thuật tái cân bằng dữ liệu trên tập huấn luyện: dùng kỹ thuật SMOTE tạo mẫu tổng hợp cho lớp thiểu số bằng cách nội suy giữa các điểm lân cận, và kỹ thuật SMOTE+ENN kết hợp thêm bước làm sạch biên quyết định nhằm loại bỏ các mẫu nhiễu trong vùng chồng lấp giữa hai lớp. Cả hai kỹ thuật chỉ được áp dụng trên tập huấn luyện để tránh việc dữ liệu được cho trước kết quả, hay còn gọi là data leakage. Chi tiết được trình bày tại mục 2.6.

\subsection{Tính phi tuyến}

Tính phi tuyến là đặc điểm của dữ liệu trong đó mối quan hệ giữa các đặc trưng đầu vào và nhãn lớp không thể biểu diễn bằng một hàm tuyến tính đơn giản. Trong thực tế lâm sàng, đây là quy luật phổ biến hơn là ngoại lệ, chẳng hạn như tác động của liều lượng thuốc lên hiệu quả điều trị, mối liên hệ giữa chỉ số BMI và nguy cơ mắc bệnh tim mạch, hay sự tương tác giữa tuổi tác và các chỉ số sinh hóa, tất cả đều mang tính phi tuyến rõ rệt. Trong dữ liệu biểu hiện gen, các gen không hoạt động độc lập mà tương tác với nhau theo các mạng lưới điều hòa phức tạp, tạo ra cấu trúc phi tuyến đặc trưng trong không gian đặc trưng. Việc này khiến các mô hình tuyến tính như Logistic Regression không đủ khả năng nắm bắt các mối quan hệ phức tạp này, dẫn đến sai lệch hệ thống trong dự đoán và hiệu năng thấp hơn so với các mô hình phi tuyến. Tuy nhiên, tính phi tuyến không phải lúc nào cũng hiển thị rõ ràng từ dữ liệu thô, cần có các phương pháp phân tích và trực quan hóa phù hợp để xác nhận sự tồn tại của nó trước khi lựa chọn mô hình.\\

Trong phạm vi khóa luận, nhóm khảo sát tính phi tuyến trên cả hai bộ dữ liệu thông qua ba hướng bổ trợ nhau. Thứ nhất, trực quan hóa t-SNE và UMAP để quan sát cấu trúc phân tách giữa các lớp trong không gian 2D phi tuyến. Thứ hai, so sánh learning curve của Logistic Regression với các mô hình phi tuyến (Random Forest, XGBoost) — nếu Logistic Regression đạt trần hiệu năng thấp hơn rõ rệt, đây là bằng chứng thực nghiệm xác nhận dữ liệu có cấu trúc phi tuyến mà mô hình tuyến tính không thể nắm bắt. Thứ ba, phân tích feature importance của Random Forest và XGBoost để xác định các đặc trưng đóng vai trò phân biệt chính, phản ánh cấu trúc phi tuyến có tổ chức trong dữ liệu.

\section{Kỹ thuật trực quan hóa của dữ liệu}
Đối với các dữ liệu số chiều cao, trực quan hóa trực tiếp trong không gian gốc là không khả thi do con người chỉ có thể nhận thức trực quan ở tối đa 3 chiều. Các kỹ thuật giảm chiều tuyến tính như PCA có thể chiếu dữ liệu xuống 2D nhưng chỉ bảo toàn phương sai tuyến tính, do đó dễ bỏ sót các cấu trúc phi tuyến phức tạp vốn đặc trưng của dữ liệu y tế. Vì vậy, đề tài sử dụng hai kỹ thuật chiếu phi tuyến là t-SNE và UMAP để quan sát cấu trúc phân bố dữ liệu trong không gian 2D — nhằm phân tích khả năng phân tách giữa các lớp và thu thập bằng chứng ban đầu về tính phi tuyến, không dùng để huấn luyện mô hình.

\subsection{t-SNE (t-Distributed Stochastic Neighbor Embedding)}
\subsection{UMAP (Uniform Manifold Approximation and Projection)} 

\section{Kỹ thuật lựa chọn đặc trưng}
\label{sec:feature_section}
\subsection{Lựa chọn đặc trưng dựa trên thống kê}
Filter methods đánh giá và lựa chọn đặc trưng dựa hoàn toàn trên các tiêu chí thống kê của dữ liệu mà không sử dụng mô hình học máy trong quá trình lựa chọn. Ưu điểm của nhóm phương pháp này là chi phí tính toán thấp, tốc độ xử lý nhanh và ít bị ảnh hưởng bởi hiện tượng overfitting.

Trong nghiên cứu này, hai bước lọc được áp dụng tuần tự gồm Variance Threshold và SelectKBest.

\paragraph{Variance Threshold}

Variance Threshold loại bỏ các đặc trưng có phương sai thấp hơn một ngưỡng định trước $\theta$. Phương sai của đặc trưng $j$ được tính theo:

\[
\mathrm{Var}(x_j)
=
\frac{1}{n}
\sum_{i=1}^{n}
(x_{ij}-\bar{x}_j)^2
\]

trong đó:

\begin{itemize}
    \item $n$ là số mẫu,
    \item $x_{ij}$ là giá trị của đặc trưng $j$ tại mẫu $i$,
    \item $\bar{x}_j$ là giá trị trung bình của đặc trưng $j$.
\end{itemize}

Những đặc trưng có $\mathrm{Var}(x_j) < \theta$ sẽ bị loại bỏ vì mức độ biến động quá thấp giữa các mẫu, từ đó khó đóng góp vào khả năng phân biệt lớp. Đây được xem là bước lọc thô ban đầu giúp giảm nhanh số lượng đặc trưng trước khi áp dụng các phương pháp tinh lọc hơn.

\paragraph{SelectKBest (ANOVA F-test)}

Sau bước Variance Threshold, phương pháp SelectKBest được sử dụng để lựa chọn $K$ đặc trưng có điểm thống kê cao nhất.

Đối với bài toán phân loại đa lớp, ANOVA F-test được sử dụng để đo mức độ khác biệt giữa các lớp. F-statistic của đặc trưng $j$ được tính như sau:

\[
F_j
=
\frac{MS_{\mathrm{between}}}
     {MS_{\mathrm{within}}}
=
\frac{
\displaystyle
\sum_{k=1}^{K}
n_k
(\bar{x}_{kj}-\bar{x}_j)^2/(K-1)
}
{
\displaystyle
\sum_{k=1}^{K}
\sum_{i\in C_k}
(x_{ij}-\bar{x}_{kj})^2/(n-K)
}
\]

trong đó:

\begin{itemize}
    \item $K$ là số lớp,
    \item $n_k$ là số mẫu thuộc lớp $k$,
    \item $\bar{x}_{kj}$ là giá trị trung bình của đặc trưng $j$ trong lớp $k$,
    \item $\bar{x}_j$ là giá trị trung bình của đặc trưng $j$ trên toàn bộ tập dữ liệu,
    \item $n$ là tổng số mẫu.
\end{itemize}

Giá trị $F_j$ càng lớn cho thấy phương sai giữa các lớp lớn hơn đáng kể so với phương sai bên trong từng lớp, đồng nghĩa với việc đặc trưng đó có khả năng phân biệt tốt hơn.

Số lượng đặc trưng được giữ lại ($K$) được xác định thông qua quá trình cross-validation.

Pipeline filter được xây dựng theo trình tự:

\[
\text{Variance Threshold}
\rightarrow
\text{SelectKBest}
\]

Kết quả thu được là một tập đặc trưng có kích thước nhỏ gọn, dễ diễn giải và có thời gian xử lý nhanh.

\subsubsection{Wrapper Methods -- Recursive Feature Elimination (RFE)}

Khác với filter methods, wrapper methods đánh giá trực tiếp các tập con đặc trưng dựa trên hiệu năng thực tế của một mô hình phân loại. Trong nghiên cứu này, phương pháp Recursive Feature Elimination (RFE) được sử dụng.

Nguyên lý hoạt động của RFE như sau:

\begin{enumerate}
    \item Huấn luyện mô hình trên toàn bộ tập đặc trưng.
    \item Tính mức độ quan trọng (feature importance) của từng đặc trưng.
    \item Loại bỏ đặc trưng có mức độ quan trọng thấp nhất.
    \item Lặp lại quá trình cho đến khi đạt số lượng đặc trưng mong muốn.
\end{enumerate}

Random Forest được sử dụng làm mô hình nền (base estimator). Mức độ quan trọng của đặc trưng được xác định thông qua Mean Decrease in Impurity (MDI):

\[
\mathrm{Importance}(x_j)
=
\frac{1}{B}
\sum_{b=1}^{B}
\sum_{t\in T_b}
\Delta i(t,x_j)
\]

trong đó:

\begin{itemize}
    \item $B$ là số cây trong Random Forest,
    \item $T_b$ là tập các nút trong cây thứ $b$,
    \item $\Delta i(t,x_j)$ là mức giảm impurity tại nút $t$ do đặc trưng $x_j$ tạo ra.
\end{itemize}

\paragraph{RFECV}

Để tự động xác định số lượng đặc trưng tối ưu, nghiên cứu sử dụng Recursive Feature Elimination with Cross-Validation (RFECV).

RFECV kết hợp RFE với Stratified K-Fold Cross-Validation. Ở mỗi kích thước tập đặc trưng, mô hình được đánh giá bằng cross-validation và điểm F1-score trung bình được sử dụng để lựa chọn số lượng đặc trưng tối ưu.

Ưu điểm của RFECV là có khả năng:

\begin{itemize}
    \item Xem xét tương tác giữa các đặc trưng.
    \item Tự động xác định số lượng đặc trưng phù hợp.
    \item Tối ưu trực tiếp cho mô hình phân loại mục tiêu.
\end{itemize}

Tuy nhiên, chi phí tính toán của RFECV cao hơn đáng kể so với filter methods do phải huấn luyện mô hình nhiều lần. Để giảm chi phí tính toán trên bộ dữ liệu BCGE có số chiều rất lớn, RFECV được thực hiện trên tập đặc trưng đã qua bước Variance Threshold thay vì trên toàn bộ tập đặc trưng gốc.

\subsubsection{So sánh hai nhóm phương pháp}

Sau khi áp dụng cả hai nhóm phương pháp, kết quả feature selection được so sánh dựa trên ba tiêu chí:

\begin{itemize}
    \item Số lượng đặc trưng được chọn.
    \item Mức độ trùng lặp giữa hai tập đặc trưng.
    \item Hiệu năng phân loại đạt được trên tập kiểm tra.
\end{itemize}

Bảng~\ref{tab:feature_selection_comparison} trình bày sự khác biệt chính giữa hai nhóm phương pháp.


Mục tiêu của việc so sánh là xác định phương pháp lựa chọn đặc trưng phù hợp hơn đối với dữ liệu gene expression có số chiều cao. Kết quả thực nghiệm chi tiết sẽ được trình bày trong Chương 4.

\section{Các mô hình phân loại}

\subsection{Logistic Regression}
\subsubsection{Khái niệm}
% ============================================================
 
Hồi quy logistic là một trong những thuật toán học máy có giám sát được sử dụng rộng rãi nhất trong lĩnh vực phân loại dữ liệu. Mặc dù mang tên hồi quy, mô hình này về bản chất là một phương pháp phân loại, được thiết kế để dự đoán xác suất thuộc về một lớp nhất định của biến mục tiêu rời rạc. Thuật toán này lần đầu tiên được giới thiệu trong lĩnh vực thống kê bởi nhà thống kê học David Cox vào năm 1958 và sau đó được tích hợp rộng rãi vào các hệ thống học máy hiện đại. Về mặt toán học, hồi quy logistic là một mô hình tuyến tính tổng quát, trong đó hàm liên kết là hàm logit. Mô hình này ánh xạ một tổ hợp tuyến tính của các biến đầu vào sang một giá trị xác suất nằm trong khoảng $(0,\,1)$, nhờ đó cho phép đưa ra quyết định phân lớp dựa trên ngưỡng xác suất được xác định trước. \\

Dạng tổng quát của hàm tuyến tính tổng quát được biểu diễn như sau:
\begin{equation}
  z = g(\mu) = \mathbf{w}^{T}\mathbf{x} + b
  \label{eq:glm}
\end{equation}

trong đó $\mu=E(Y|X)$ là giá trị kỳ vọng của biến phản hồi, $g(\cdot)$ là hàm liên kết, $\mathbf{x}$ là vector đặc trưng đầu vào, $\mathbf{w}$ là vector trọng số và $b$ là hệ số chệch.\\

Thấy rằng, để dự đoán xác suất thuộc lớp dương, giá trị đầu ra $z$ có thể nhận bất kỳ giá trị thực nào trong khoảng $(-\infty,+\infty)$. Điều này không phù hợp với yêu cầu của bài toán phân loại, bởi xác suất phải nằm trong khoảng $[0,1]$. Ví dụ, mô hình tuyến tính có thể cho kết quả $z=2.5$ hoặc $z=-1.3$, những giá trị không thể được diễn giải như xác suất. Mô hình logistic được xây dựng nhằm giải quyết các bài toán phân loại nhị phân, tức là bài toán trong đó biến phụ thuộc chỉ nhận một trong hai giá trị, thường được ký hiệu là $0$ (âm tính) và $1$ (dương tính). Việc quyết định lớp được thực hiện dựa trên một ngưỡng xác suất nhất định (thường là 0.5). Chẳng hạn, trong lĩnh vực y tế, mô hình có thể được sử dụng để dự đoán liệu một bệnh nhân có mắc bệnh tiểu đường hay không dựa trên các chỉ số sinh học. Hay trong lĩnh vực thể thao, mô hình có thể được sử dụng để dự đoán xác suất ghi bàn của một cú sút (vào hoặc không vào). Ngoài ra, trong tài chính, mô hình giúp xác định khả năng khách hàng có vỡ nợ hay không. Tất cả các quyết định này đều dựa trên giá trị xác suất được dự đoán bởi mô hình thông qua các đặc trưng đầu vào tương ứng.\\

Ngoài bài toán nhị phân, hồi quy logistic có thể được mở rộng sang bài toán phân loại đa lớp, tức là bài toán trong đó biến mục tiêu $y$ có thể nhận một trong $K > 2$ giá trị rời rạc thay vì chỉ hai giá trị (0 hoặc 1) như thông thường. Bằng cách tiếp cận theo các hướng phổ biến như phân loại nhiều lớp với mỗi lớp là một bài toán nhị phân độc lập (One-vs-Rest) hoặc xây dựng các bộ phân loại trên từng cặp lớp (One-vs-One), hoặc mở rộng trực tiếp hàm sigmoid thành hàm softmax, hồi quy logistic có thể được áp dụng hiệu quả cho các bài toán phân loại đa lớp. Trong phạm vi luận văn này, đối với bộ dữ liệu Breast Cancer Gene Expression gồm sáu lớp phân loại, mô hình được triển khai dưới dạng hồi quy logistic đa lớp, trong đó hàm softmax được sử dụng để ước lượng xác suất thuộc về từng lớp. Lớp có xác suất dự đoán cao nhất sẽ được lựa chọn làm kết quả phân loại cuối cùng. \\

\subsubsection{Hàm Sigmoid}

Như đề cập ở trên, hàm sigmoid (còn gọi là hàm logistic) có thể được xem như là "trái tim" của mô hình hồi quy logistic. Hàm này được định nghĩa như sau:
 
\begin{equation}
  \sigma(z) = \frac{1}{1 + e^{-z}}
  \label{eq:sigmoid}
\end{equation}
 
\noindent trong đó $z$ là giá trị đầu vào thực bất kỳ và $e$ là cơ số của logarithm tự nhiên. Tính chất của hàm sigmoid là: (1) Miền giá trị của hàm nằm hoàn toàn trong khoảng $(0,\,1)$, điều này khiến đầu ra của hàm có thể được diễn giải một cách tự nhiên như xác suất; (2) Hàm có dạng chữ S đơn điệu tăng, nghĩa là khi $z \to +\infty$ thì $\sigma(z) \to 1$, và khi $z \to -\infty$ thì $\sigma(z) \to 0$; (3) Tại $z = 0$, hàm cho giá trị $\sigma(0) = 0{,}5$, tương ứng với trạng thái không chắc chắn hoàn toàn giữa hai lớp. \\

Khi đạo hàm, hàm có dạng:\\

\begin{equation}
  \sigma'(z) = \sigma(z)\bigl(1 - \sigma(z)\bigr)
  \label{eq:sigmoid_deriv}
\end{equation}
 
\noindent Khi đó, hàm có dạng thuận lợi trong tính toán, giúp đơn giản hóa đáng kể quá trình tối ưu hóa tham số bằng phương pháp lặp cập nhật tham số theo hướng giảm nhanh nhất của hàm mất mát (được đề cập tại mục \ref{sec:gradient}). Về mặt ý nghĩa thực tiễn, hàm sigmoid đóng vai trò là cầu nối chuyển đổi một đầu ra tuyến tính không bị giới hạn thành một xác suất có ý nghĩa thống kê, khắc phục nhược điểm chính của hồi quy tuyến tính khi áp dụng trực tiếp cho bài toán phân loại.\\

\subsubsection{Nguyên lý hoạt động}
% ============================================================
 
Hồi quy logistic hoạt động dựa trên ý tưởng cốt lõi là ước lượng xác suất có điều kiện $P(y = 1 \mid \mathbf{x})$, tức là xác suất để quan sát $\mathbf{x}$ thuộc về lớp dương tính. Thay vì dự đoán trực tiếp nhãn lớp, mô hình tính toán một giá trị xác suất liên tục, sau đó áp dụng một ngưỡng quyết định (thường là $0{,}5$) để gán nhãn cuối cùng. Quá trình hoạt động của mô hình có thể được mô tả qua ba bước chính. Trước tiên, mô hình tính toán tổ hợp tuyến tính $z = \mathbf{w}^{T}\mathbf{x} + b$ từ các đặc trưng đầu vào. Tiếp theo, giá trị $z$ được đưa qua hàm sigmoid để chuyển đổi thành xác suất $\hat{p} = \sigma(z)$. Cuối cùng, dựa trên $\hat{p}$, mô hình đưa ra quyết định phân lớp theo ngưỡng đã định.\\
 
Mô hình hồi quy logistic được định nghĩa hoàn chỉnh thông qua công thức sau:
 
\begin{equation}
  \hat{p} = P(y = 1 \mid \mathbf{x})
           = \sigma\!\left(\mathbf{w}^{T}\mathbf{x} + b\right)
           = \frac{1}{1 + e^{-(\mathbf{w}^{T}\mathbf{x} + b)}}
  \label{eq:logreg}
\end{equation}
 
\noindent Trong đó:
 
\begin{itemize}
  \item $\mathbf{x} = (x_1, x_2, \ldots, x_n)^{T} \in \mathbb{R}^{n}$: Vector đặc trưng đầu vào gồm $n$ biến độc lập, mô tả các thuộc tính của mẫu dữ liệu cần phân loại.
 
  \item $\mathbf{w} = (w_1, w_2, \ldots, w_n)^{T} \in \mathbb{R}^{n}$: Vector trọng số - các tham số học được từ dữ liệu huấn luyện. Mỗi $w_i$ thể hiện mức độ ảnh hưởng của đặc trưng $x_i$ đến xác suất dự đoán. Trọng số dương đồng nghĩa với việc đặc trưng tương ứng làm tăng xác suất thuộc lớp $1$, trong khi trọng số âm thể hiện tác động ngược lại.
 
  \item $b \in \mathbb{R}$: Hệ số chệch - một hằng số cho phép mô hình linh hoạt trong việc dịch chuyển đường quyết định khỏi gốc tọa độ, giúp mô hình khớp tốt hơn với dữ liệu thực tế.
 
  \item $z = \mathbf{w}^{T}\mathbf{x} + b$: Giá trị tổng hợp tuyến tính, còn được gọi là \textit{log-odds} hay \textit{logit}, biểu thị logarithm của tỉ lệ xác suất:
  \[
    z = \log\!\left(\frac{\hat{p}}{1 - \hat{p}}\right)
  \]
 
  \item $\hat{p}$: Xác suất ước lượng rằng mẫu $\mathbf{x}$ thuộc về lớp dương tính ($y = 1$).
\end{itemize}
 
Sau khi tính được xác suất $\hat{p}$, mô hình đưa ra quyết định phân lớp dựa trên ngưỡng quyết định $\tau$ (thường là $0{,}5$) theo quy tắc:
 
\begin{equation}
  \hat{y} =
  \begin{cases}
    1 & \text{nếu } \hat{p} \geq \tau \\
    0 & \text{nếu } \hat{p} < \tau
  \end{cases}
  \label{eq:decision}
\end{equation}
 
Ngưỡng $\tau = 0{,}5$ là lựa chọn mặc định và hợp lý trong điều kiện dữ liệu cân bằng. Tuy nhiên, trong thực tế, khi bài toán yêu cầu độ nhạy cao - ví dụ: phát hiện ung thư, trong đó bỏ sót dương tính thật có chi phí rất cao - ngưỡng $\tau$ có thể được điều chỉnh xuống thấp hơn để tăng chỉ số Recall nhưng chấp nhận giảm độ chính xác (chỉ số Precision). Việc lựa chọn ngưỡng phù hợp thường được hỗ trợ bởi đường cong ROC và chỉ số AUC (chỉ số ROC-AUC được đề cập trong mục \ref{sec:matcanbang}).
 
% ============================================================
\subsubsection{Hàm mất mát và phương pháp tối ưu}
\label{sec:gradient}

Như đã trình bày, mô hình cần xác định một tiêu chí đo lường mức độ sai lệch giữa xác suất mô hình dự đoán và nhãn thực tế trong tập dữ liệu. Hồi quy logistic sử dụng hàm mất mát Log-Loss, được dẫn xuất từ nguyên lý ước lượng hợp lý cực đại. Với tập huấn luyện gồm $m$ mẫu, hàm mất mát được định nghĩa như sau:

\begin{equation}
  \mathcal{L}(\mathbf{w}, b)
  = -\frac{1}{m} \sum_{i=1}^{m}
    \Bigl[
      y^{(i)} \log \hat{p}^{(i)}
      + \bigl(1 - y^{(i)}\bigr) \log \bigl(1 - \hat{p}^{(i)}\bigr)
    \Bigr]
  \label{eq:loss}
\end{equation}
 
\noindent trong đó $m$ là số lượng mẫu huấn luyện, $y^{(i)}$ là nhãn thực tế và $\hat{p}^{(i)}$ là xác suất dự đoán của mẫu thứ $i$.

Về mặt trực quan, hàm này phạt nặng những trường hợp mô hình dự đoán sai với mức độ tự tin cao, chẳng hạn, dự đoán xác suất gần 0 cho một mẫu thực tế thuộc lớp dương, hoặc ngược lại. Hàm mất mát này có tính chất lồi đối với không gian tham số, đảm bảo rằng mọi nghiệm cực bộ cũng là nghiệm toàn cục — một tính chất quan trọng giúp quá trình tối ưu hóa hội tụ ổn định. Để cực tiểu hóa $\mathcal{L}$, phương pháp phổ biến nhất là gradient descent (tạm dịch là thuật toán hạ gradient), trong đó các tham số được cập nhật lặp đi lặp lại theo hướng ngược chiều gradient:

\begin{equation}
  \mathbf{w} \leftarrow \mathbf{w} - \eta \cdot \frac{\partial \mathcal{L}}{\partial \mathbf{w}},
  \qquad
  b \leftarrow b - \eta \cdot \frac{\partial \mathcal{L}}{\partial b}
  \label{eq:update}
\end{equation}
 
\noindent trong đó $\eta > 0$ là tốc độ học - một siêu tham số kiểm soát bước dịch chuyển trong không gian tham số.\\
 
Ngoài ra, để tránh hiện tượng mô hình học vẹt, người ta thường bổ sung thêm các thành phần chính quy hóa (Regularization) vào hàm mất mát. Trong phạm vi khóa luận này, chính quy hóa L2 (\textit{Ridge}) được sử dụng:
 
\begin{equation}
  \mathcal{L}_{\mathrm{reg}}
  = \mathcal{L} + \lambda \|\mathbf{w}\|_{2}^{2}
  \label{eq:l2}
\end{equation}
 
\noindent trong đó $\lambda > 0$ là hệ số chính quy hóa, kiểm soát mức độ phạt đối với các trọng số lớn. \\

% ============================================================
\subsubsection{Ưu điểm} 
Hồi quy logistic có một số đặc điểm thuận lợi khiến mô hình này được sử dụng rộng rãi trong thực tế. Thứ nhất, mô hình có tính diễn giải cao: các hệ số $w_i$ phản ánh trực tiếp tác động của từng đặc trưng đến log-odds của kết quả, cho phép phân tích ý nghĩa của từng biến đầu vào một cách tường minh. Nhờ đó, mô hình được ưu tiên trong các lĩnh vực đòi hỏi tính minh bạch như y tế - nơi cần giải thích nguy cơ mắc bệnh dựa trên chỉ số lâm sàng - hay tài chính, nơi khả năng lý giải kết quả cho các bên liên quan là yêu cầu bắt buộc trong các hệ thống chấm điểm tín dụng và phát hiện gian lận.

Thứ hai, mô hình có độ phức tạp tính toán thấp, phù hợp với các tập dữ liệu lớn. Do hàm mất mát có tính lồi, quá trình huấn luyện hội tụ ổn định và không phụ thuộc vào giá trị khởi tạo tham số. Ngoài ra, đầu ra dạng xác suất cho phép điều chỉnh ngưỡng quyết định linh hoạt, làm cho mô hình phù hợp với nhiều bài toán thực tế như dự đoán hành vi người dùng (dự đoán tỷ lệ khách hàng rời bỏ ) hay phân loại văn bản (phát hiện tin rác, phân tích cảm xúc).
\subsubsection{Hạn chế} Hồi quy logistic cũng tồn tại một số hạn chế cần lưu ý. Hạn chế cơ bản nhất là giả định tuyến tính: mô hình chỉ có khả năng học các ranh giới quyết định tuyến tính trong không gian đặc trưng, do đó hiệu suất có thể bị giới hạn với các bài toán có cấu trúc phi tuyến phức tạp. Ngoài ra, mô hình nhạy cảm với đa cộng tuyến — khi các biến đầu vào có tương quan cao với nhau, ước lượng tham số có thể trở nên không ổn định. Cuối cùng, trong bối cảnh dữ liệu mất cân bằng lớp, hiệu suất mô hình có thể suy giảm đáng kể nếu không áp dụng các kỹ thuật hỗ trợ phù hợp.\\
% ============================================================

\subsection{Mô hình SVM (Support Vector Machine)}
% ============================================================
\subsubsection{Khái niệm}
% ============================================================

Support Vector Machine (SVM) là một thuật toán học máy có giám sát được phát triển bởi Vladimir Vapnik và Alexey Chervonenkis vào những năm 1960, sau đó được hoàn thiện và phổ biến rộng rãi thông qua công trình của Cortes và Vapnik năm 1995. SVM được xây dựng trên nền tảng lý thuyết tối thiểu hóa rủi ro có cấu trúc, nhằm tìm kiếm một mô hình không chỉ khớp tốt với dữ liệu huấn luyện mà còn có khả năng tổng quát hóa cao trên dữ liệu chưa thấy. Về bản chất, SVM là một bộ phân loại tuyến tính tìm kiếm siêu phẳng tối ưu phân tách các lớp dữ liệu trong không gian đặc trưng, sao cho khoảng cách từ siêu phẳng đến các điểm dữ liệu gần nhất của mỗi lớp là lớn nhất có thể. Thông qua kỹ thuật hàm nhân, SVM có thể mở rộng sang các bài toán phi tuyến mà không cần tường minh tính toán trong không gian chiều cao.\\

Cũng như hồi quy logistic, SVM được thiết kế chủ yếu cho các bài toán phân loại nhị phân, trong đó mục tiêu là phân chia không gian đặc trưng thành hai vùng tương ứng với hai nhãn lớp. Thông qua các phương pháp mở rộng phổ biến như phân loại nhiều lớp với mỗi lớp là một bài toán nhị phân độc lập (One-vs-Rest) hoặc xây dựng các bộ phân loại trên từng cặp lớp (One-vs-One), SVM có thể xử lý hiệu quả các bài toán phân loại đa lớp. Ngoài ra, biến thể hồi quy vector hỗ trợ (Support Vector Regression - SVR) cho phép áp dụng nguyên lý của SVM vào các bài toán hồi quy, mở rộng đáng kể phạm vi ứng dụng của mô hình.\\

% ============================================================
\subsubsection{Siêu phẳng, margin và support vectors}
% ============================================================

Trong không gian $n$ chiều $\mathbb{R}^n$, một siêu phẳng là tập hợp các điểm $\mathbf{x}$ thỏa mãn phương trình tuyến tính:\\

\begin{equation}
    \mathbf{w}^T \mathbf{x} + b = 0 \tag{2.9}
\end{equation}

\noindent trong đó $\mathbf{w} \in \mathbb{R}^n$ là vector pháp tuyến xác định hướng của siêu phẳng và $b \in \mathbb{R}$ là hệ số chệch xác định vị trí của siêu phẳng trong không gian. Siêu phẳng này chia không gian thành hai nửa tương ứng với hai lớp phân loại: lớp dương tính ($y = +1$) khi $\mathbf{w}^T\mathbf{x} + b > 0$ và lớp âm tính ($y = -1$) khi $\mathbf{w}^T\mathbf{x} + b < 0$.\\

Khoảng cách từ một điểm $\mathbf{x}_i$ đến siêu phẳng được xác định bởi:\\

\begin{equation}
    d_i = \frac{|\mathbf{w}^T\mathbf{x}_i + b|}{\|\mathbf{w}\|} \tag{2.10}
\end{equation}

Biên độ (\textit{margin}) của siêu phẳng được định nghĩa là tổng khoảng cách từ siêu phẳng đến điểm gần nhất của mỗi lớp. Khi chuẩn hóa sao cho $|\mathbf{w}^T\mathbf{x}_i + b| = 1$ với các điểm gần nhất, biên độ này bằng $\dfrac{2}{\|\mathbf{w}\|}$. Mục tiêu của SVM là tìm siêu phẳng có biên độ lớn nhất, vì biên độ càng lớn thì mô hình càng ít nhạy cảm với nhiễu và có khả năng tổng quát hóa tốt hơn.\\

Các vector hỗ trợ (\textit{support vectors}) là những điểm dữ liệu nằm trên biên giới của biên độ, tức là những điểm thỏa mãn $|\mathbf{w}^T\mathbf{x}_i + b| = 1$. Đây là các điểm duy nhất quyết định vị trí và hướng của siêu phẳng tối ưu --- nếu loại bỏ các điểm này, siêu phẳng sẽ thay đổi, trong khi loại bỏ các điểm còn lại thì không.\\
% ============================================================

\subsubsection{Nguyên lý hoạt động}
% ============================================================

Ý tưởng cốt lõi của SVM là tìm một ranh giới quyết định phân tách hai lớp dữ liệu sao cho khoảng cách từ ranh giới đó đến các điểm dữ liệu gần nhất của mỗi lớp được tối đa hóa. Bằng cách tìm siêu phẳng tối đa hóa biên độ $\dfrac{2}{\|\mathbf{w}\|}$ (tương đương với việc tối thiểu hóa $\|\mathbf{w}\|^2$), đồng thời đảm bảo tất cả các điểm dữ liệu được phân loại đúng. Có thể nói, lập luận cơ bản đằng sau nguyên tắc này là: trong số tất cả các siêu phẳng có thể phân tách hai lớp, siêu phẳng có biên độ lớn nhất cung cấp sự phân tách ``an toàn'' nhất, làm giảm thiểu rủi ro phân loại sai trên dữ liệu mới. Nguyên lý này xuất phát từ lập luận rằng một ranh giới quyết định có biên độ lớn hơn sẽ có độ tin cậy cao hơn và ít nhạy cảm hơn với nhiễu trong dữ liệu, từ đó dẫn đến khả năng tổng quát hóa tốt hơn trên dữ liệu mới.\\

Vậy nên, quá trình huấn luyện của SVM về bản chất là một bài toán tối ưu hóa lồi, trong đó mô hình tìm kiếm các bộ tham số xác định siêu phẳng tối ưu thông qua phương pháp gradient descent - phương pháp lặp cập nhật tham số theo chiều giảm nhanh nhất của hàm mục tiêu. Vì SVM chỉ phụ thuộc vào một tập con nhỏ là các điểm dữ liệu huấn luyện, hay còn gọi là các vector hỗ trợ để xác định ranh giới quyết định, thay vì toàn bộ tập dữ liệu. Tính chất này mang lại cho SVM sự hiệu quả đặc biệt về mặt bộ nhớ và độ bền trước các điểm dữ liệu ngoại lệ nằm xa ranh giới.


\subsubsection{Soft-Margin SVM}

Trong thực tế, dữ liệu hiếm khi có thể được phân tách hoàn toàn bằng một siêu phẳng tuyến tính - đặc biệt khi dữ liệu có nhiễu hoặc các lớp có vùng chồng lấp nhau. Để xử lý trường hợp này, SVM được mở rộng thành biến thể soft-margin bằng cách cho phép một số điểm dữ liệu vi phạm ràng buộc phân tách, thay vì yêu cầu tất cả các điểm phải nằm đúng phía và cách xa siêu phẳng ít nhất một khoảng bằng $\dfrac{1}{\|\mathbf{w}\|}$.\\

Cụ thể, với mỗi điểm dữ liệu $\mathbf{x}_i$, một biến bù $\xi_i \geq 0$ được đưa vào để đo mức độ vi phạm: $\xi_i = 0$ nghĩa là điểm nằm đúng phía và đủ xa siêu phẳng, $0 < \xi_i \leq 1$ nghĩa là điểm nằm đúng phía nhưng quá gần siêu phẳng, còn $\xi_i > 1$ nghĩa là điểm bị phân loại sai. Bài toán tối ưu hóa lúc này trở thành:\\

\begin{equation}
    \min_{\mathbf{w}, b, \xi} \quad \frac{1}{2}\|\mathbf{w}\|^2 + C\sum_{i=1}^{m} \xi_i \tag{2.11}
\end{equation}

\begin{equation}
    \text{ràng buộc} \quad y_i(\mathbf{w}^T\mathbf{x}_i + b) \geq 1 - \xi_i, \quad \xi_i \geq 0, \quad \forall i \tag{2.12}
\end{equation}

Tham số $C > 0$ là siêu tham số điều tiết, kiểm soát sự đánh đổi giữa việc tối đa hóa biên độ và chấp nhận vi phạm phân loại. Khi $C$ lớn, mô hình bị phạt nặng hơn với mỗi vi phạm, do đó ưu tiên phân loại đúng tất cả các điểm huấn luyện nhưng có thể dẫn đến biên độ nhỏ hơn và nguy cơ quá khớp cao hơn. Ngược lại, khi $C$ nhỏ, mô hình chấp nhận nhiều vi phạm hơn để đổi lấy biên độ lớn hơn, giúp mô hình tổng quát hóa tốt hơn trên dữ liệu mới.\\

\subsubsection{Kernel SVM}

Nhiều bài toán phân loại có ranh giới quyết định không thể biểu diễn bằng một siêu phẳng tuyến tính trong không gian đặc trưng gốc. Để giải quyết hạn chế này, SVM được mở rộng thông qua kỹ thuật hàm nhân (kernel), cho phép mô hình học các ranh giới quyết định phi tuyến mà không cần tường minh ánh xạ dữ liệu sang không gian chiều cao.\\

Ý tưởng cốt lõi của kỹ thuật hàm nhân xuất phát từ quan sát rằng trong bài toán đối ngẫu, dữ liệu chỉ xuất hiện dưới dạng tích vô hướng $\mathbf{x}_i^T\mathbf{x}_j$. Thay vì tường minh tính toán $\phi(\mathbf{x})$, hàm ánh xạ dữ liệu sang không gian chiều cao ta có thể thay thế tích vô hướng bằng một hàm nhân $K(\mathbf{x}_i, \mathbf{x}_j) = \phi(\mathbf{x}_i)^T\phi(\mathbf{x}_j)$, cho phép mô hình ngầm định hoạt động trong không gian chiều cao mà không cần tính toán $\phi(\mathbf{x})$ một cách tường minh. Lợi thế của kernel SVM là khả năng học các cấu trúc phi tuyến phức tạp với chi phí tính toán hợp lý.\\

Trong phạm vi khóa luận này, hàm nhân được sử dụng là hàm nhân cơ sở hướng tâm (Radial Basis Function - RBF) sẽ được sử dụng, đây cũng là hàm nhân phổ biến nhất trong thực tế:\\

\begin{equation}
    K(\mathbf{x}_i, \mathbf{x}_j) = \exp\left(-\gamma\|\mathbf{x}_i - \mathbf{x}_j\|^2\right), \quad \gamma > 0 \tag{2.13}
\end{equation}

Hàm nhân RBF đo lường sự tương đồng giữa hai điểm dữ liệu dựa trên khoảng cách Euclidean giữa chúng: hai điểm càng gần nhau thì giá trị hàm nhân càng tiệm cận 1 và càng xa nhau thì càng tiệm cận 0. Tham số $\gamma > 0$ kiểm soát phạm vi ảnh hưởng của mỗi điểm dữ liệu: $\gamma$ lớn làm cho mô hình chỉ chú ý đến các điểm lân cận rất gần, dẫn đến ranh giới quyết định cục bộ và phức tạp hơn; trong khi $\gamma$ nhỏ tạo ra ranh giới mượt mà hơn và có khả năng tổng quát hóa tốt hơn trên dữ liệu mới.\\

% ============================================================
\subsubsection*{Hàm mất mát và phương pháp tối ưu}

SVM sử dụng hàm mất mát Hinge Loss, được định nghĩa cho từng mẫu như sau:

\begin{equation}
    \ell_i = \max\left(0, 1 - y_i(\mathbf{w}^T\mathbf{x}_i + b)\right) \tag{2.14}
\end{equation}

\noindent trong đó $y_i \in \{-1, +1\}$ là nhãn thực tế và $\mathbf{w}^T\mathbf{x}_i + b$ là giá trị dự đoán của mô hình trước khi phân loại. Tích $y_i(\mathbf{w}^T\mathbf{x}_i + b)$ đo mức độ tự tin của mô hình: giá trị dương lớn nghĩa là phân loại đúng với độ tự tin cao, giá trị âm nghĩa là phân loại sai.\\

Hàm này có ý nghĩa trực quan: nếu điểm dữ liệu được phân loại đúng và nằm đủ xa siêu phẳng, mất mát bằng 0, mô hình không bị phạt. Ngược lại, nếu điểm nằm quá gần hoặc bị phân loại sai, mất mát tăng tuyến tính theo mức độ vi phạm. Nói cách khác, SVM không quan tâm đến việc phân loại đúng đến mức nào, mà chỉ phạt những trường hợp chưa đủ tự tin hoặc sai. Khi hàm mục tiêu tổng thể kết hợp Hinge Loss với thành phần chính quy hóa L2, nó có dạng như sau:

\begin{equation}
    \mathcal{L}(\mathbf{w}, b) = \frac{1}{2}\|\mathbf{w}\|^2 + C\sum_{i=1}^{m} \max\left(0, 1 - y_i(\mathbf{w}^T\mathbf{x}_i + b)\right) \tag{2.15}
\end{equation}

\noindent trong đó tham số $C > 0$ kiểm soát sự đánh đổi giữa biên độ lớn và mức độ chấp nhận vi phạm: $C$ lớn nghĩa là mô hình bị phạt nặng hơn khi có điểm bị phân loại sai, do đó cố gắng phân loại đúng tất cả các điểm huấn luyện; $C$ nhỏ nghĩa là mô hình chấp nhận một số vi phạm để đổi lấy biên độ lớn hơn và khả năng tổng quát hóa tốt hơn.

\subsubsection*{Ưu điểm và ứng dụng thực tế}

SVM có một số đặc điểm thuận lợi nổi bật trong các bài toán phân loại. Thứ nhất, nhờ nguyên tắc tối đa hóa biên độ, mô hình có khả năng tổng quát hóa tốt và giảm thiểu rủi ro quá khớp ngay cả trong không gian chiều cao. Thứ hai, thông qua kỹ thuật hàm nhân, SVM có thể xử lý hiệu quả các bài toán phi tuyến phức tạp mà không bị ảnh hưởng bởi lời nguyền chiều cao như các phương pháp khác. Thứ ba, do chỉ phụ thuộc vào các vector hỗ trợ, mô hình hiệu quả về bộ nhớ và bền vững trước các điểm dữ liệu ngoại lệ nằm xa ranh giới quyết định.

Nhờ các đặc điểm trên, SVM được ứng dụng rộng rãi trong nhiều lĩnh vực đòi hỏi độ chính xác cao. Trong y sinh học, mô hình được sử dụng để phân loại dữ liệu gen, chẩn đoán bệnh từ dữ liệu y tế và phân tích tín hiệu não (EEG/fMRI). Trong phân loại văn bản và xử lý ngôn ngữ tự nhiên, SVM với kernel tuyến tính cho kết quả tốt trong không gian đặc trưng thưa chiều cao. Trong tài chính, mô hình được sử dụng cho phát hiện gian lận và dự đoán biến động thị trường.

\subsubsection*{Hạn chế}

Bên cạnh các ưu điểm trên, SVM cũng tồn tại một số hạn chế cần lưu ý. Hạn chế lớn nhất là chi phí tính toán cao khi huấn luyện trên tập dữ liệu lớn: độ phức tạp thời gian của các thuật toán giải bài toán tối ưu thường là $O(m^2)$ đến $O(m^3)$ theo số mẫu $m$, khiến SVM trở nên kém hiệu quả khi $m$ lên đến hàng triệu. Thứ hai, việc lựa chọn kernel và điều chỉnh các siêu tham số ($C$, $\gamma$) đòi hỏi phải thử nghiệm nhiều tổ hợp giá trị khác nhau và đánh giá hiệu suất từng tổ hợp, làm tăng đáng kể chi phí tính toán tổng thể. Cuối cùng, không giống như hồi quy logistic, SVM không trực tiếp cho ra xác suất mà chỉ cho ra nhãn phân loại. Để có được đầu ra dạng xác suất, cần huấn luyện thêm một mô hình phụ trợ trên đầu ra của SVM, làm tăng thêm chi phí tính toán so với các mô hình vốn đã cho ra xác suất trực tiếp.

\subsection{Mô hình Random Forest}
\subsubsection{Khái niệm}

Rừng ngẫu nhiên (Random Forest) là một thuật toán học máy có giám sát thuộc lớp các phương pháp học tập kết hợp (ensemble learning), được xây dựng bằng cách kết hợp một tập hợp lớn các cây quyết định được huấn luyện độc lập trên các tập con dữ liệu ngẫu nhiên khác nhau. Kết quả dự đoán cuối cùng được tổng hợp từ tất cả các cây thành phần thông qua cơ chế biểu quyết đa số trong bài toán phân loại, hoặc lấy trung bình trong bài toán hồi quy. Tên gọi "rừng ngẫu nhiên" phản ánh hai yếu tố cốt lõi: "rừng" biểu thị tập hợp nhiều cây quyết định, và "ngẫu nhiên" biểu thị hai nguồn ngẫu nhiên hóa được đưa vào, lấy mẫu tập dữ liệu có thay thế và lựa chọn ngẫu nhiên tập con đặc trưng tại mỗi nút phân chia.\\

Về mặt hình thức, cho tập huấn luyện $\mathcal{D} = \{(\mathbf{x}_i, y_i)\}_{i=1}^m$ gồm $m$ mẫu, Random Forest xây dựng một tập hợp $B$ cây quyết định $\{T_1, T_2, \ldots, T_B\}$, trong đó mỗi cây $T_b$ được huấn luyện trên tập bootstrap $\mathcal{D}_b$, tức là tập dữ liệu được lấy mẫu có thay thế từ $\mathcal{D}$ với cùng kích thước $m$ và sử dụng ngẫu nhiên một tập con gồm $k$ đặc trưng tại mỗi nút phân chia thay vì toàn bộ đặc trưng. Ở đây $\mathbf{x}_i \in \mathbb{R}^n$ là vector đặc trưng và $y_i$ là nhãn lớp của mẫu thứ $i$. Dự đoán cuối cùng là hàm tổng hợp trên toàn bộ $B$ cây.\\

Thuật toán Random Forest được Leo Breiman công bố chính thức vào năm 2001, tổng hợp và hệ thống hóa các ý tưởng từ nhiều công trình trước đó. Kỹ thuật lấy mẫu có thay thế để giảm phương sai của mô hình được đề xuất bởi Breiman từ năm 1996, trong khi kỹ thuật lựa chọn ngẫu nhiên tập con đặc trưng tại mỗi nút phân chia được nghiên cứu bởi Ho (1995) và Amit và Geman (1997). Breiman chứng minh rằng hiệu suất tổng quát hóa của Random Forest phụ thuộc vào độ mạnh trung bình của các cây thành phần và sự tương quan giữa chúng - đặt nền tảng lý thuyết cho nguyên tắc thiết kế: tối đa hóa độ mạnh của từng cây trong khi tối thiểu hóa sự tương quan giữa các cây.\\

Trong thực tế, Random Forest được thiết kế để giải quyết cả hai nhóm bài toán học máy có giám sát chính là phân loại (classification) và hồi quy (regression). Tuy nhiên, mục tiêu sâu xa hơn của thuật toán là khắc phục các hạn chế cơ bản của Decision Tree đơn lẻ, đặc biệt là tính không ổn định và xu hướng học vẹt, bằng việc kết hợp nhiều mô hình yếu thành một mô hình tổng hợp mạnh và ổn định hơn. Ngoài hai ứng dụng phân loại và hồi quy, Random Forest còn được sử dụng rộng rãi như một công cụ đánh giá tầm quan trọng đặc trưng (kỹ thuật lựa chọn đặc trưng được nhắc đến tại \ref{sec:feature_importance}), hỗ trợ lựa chọn đặc trưng trong các quy trình tiền xử lý dữ liệu phức tạp.\\

\subsubsection*{Nguyên lý hoạt động}


Ý tưởng cốt lõi của Random Forest là xây dựng nhiều cây quyết định độc lập, mỗi cây được huấn luyện trên một phiên bản ngẫu nhiên khác nhau của tập dữ liệu, sau đó tổng hợp kết quả từ tất cả các cây để đưa ra dự đoán cuối cùng. Sự kết hợp này giúp khắc phục hạn chế cơ bản của cây quyết định đơn lẻ, vốn có xu hướng quá khớp với dữ liệu huấn luyện bằng cách trung hòa các sai lệch ngẫu nhiên của từng cây.\\

Quá trình huấn luyện mỗi cây $T_b$ gồm hai bước ngẫu nhiên hóa. Thứ nhất, thay vì huấn luyện trên toàn bộ tập dữ liệu $\mathcal{D}$, mỗi cây được huấn luyện trên một tập con $\mathcal{D}_b$ được lấy mẫu có thay thế từ $\mathcal{D}$ với cùng kích thước $m$, tức là một số mẫu có thể được chọn nhiều lần trong khi một số khác không được chọn. Thứ hai, tại mỗi nút phân chia, thay vì xem xét toàn bộ $n$ đặc trưng, mô hình chỉ xem xét một tập con gồm $k$ đặc trưng được chọn ngẫu nhiên, thường lấy $k = \sqrt{n}$ cho bài toán phân loại. Hai nguồn ngẫu nhiên hóa này đảm bảo các cây trong rừng đủ khác biệt nhau, tránh trường hợp tất cả các cây đều mắc cùng một loại sai lầm. Sau khi huấn luyện $B$ cây, dự đoán cuối cùng cho một mẫu mới $\mathbf{x}$ được tổng hợp bằng cơ chế biểu quyết đa số:\\

\begin{equation}
    \hat{y} = \arg\max_{c} \sum_{b=1}^{B} \mathbf{1}\left[T_b(\mathbf{x}) = c\right] \tag{2.16}
\end{equation}

\noindent trong đó $T_b(\mathbf{x})$ là nhãn dự đoán của cây thứ $b$ cho mẫu $\mathbf{x}$, $c$ là nhãn lớp, và $\mathbf{1}[\cdot]$ là hàm chỉ thị nhận giá trị 1 nếu điều kiện đúng và 0 nếu sai. Nói cách khác, nhãn được dự đoán bởi nhiều cây nhất sẽ là kết quả cuối cùng.

\subsubsection*{Các siêu tham số}

Siêu tham số là các tham số được thiết lập trước khi huấn luyện mô hình và không được học từ dữ liệu, khác với các tham số như trọng số $\mathbf{w}$ được tối ưu hóa trong quá trình huấn luyện. Việc lựa chọn siêu tham số phù hợp ảnh hưởng trực tiếp đến hiệu suất và khả năng tổng quát hóa của mô hình. Trong phạm vi khóa luận, hai siêu tham số chính của Random Forest được điều chỉnh. Thứ nhất là số lượng cây $B$ (\texttt{n\_estimators}), kiểm soát độ ổn định của mô hình: $B$ càng lớn thì dự đoán càng ổn định do trung hòa được nhiều sai lệch ngẫu nhiên hơn, nhưng chi phí tính toán cũng tăng tương ứng. Thứ hai là chiều sâu tối đa của mỗi cây (\texttt{max\_depth}), kiểm soát mức độ phức tạp của từng cây thành phần: cây càng sâu thì càng có khả năng học các cấu trúc phức tạp trong dữ liệu, nhưng cũng dễ quá khớp hơn. Khi \texttt{max\_depth = None}, tức không dùng đến \texttt{max\_depth}, cây được phép phát triển đến khi tất cả các lá thuần nhất hoàn toàn.\\

\subsubsection*{Ưu điểm và ứng dụng thực tế}

Random Forest có một số đặc điểm thuận lợi khiến mô hình này được sử dụng rộng rãi. Thứ nhất, nhờ cơ chế kết hợp nhiều cây độc lập, mô hình có khả năng tổng quát hóa tốt và ít nhạy cảm với nhiễu trong dữ liệu so với cây quyết định đơn lẻ. Thứ hai, mô hình hoạt động hiệu quả trên cả dữ liệu chiều cao lẫn dữ liệu có nhiều đặc trưng không liên quan, nhờ cơ chế lựa chọn ngẫu nhiên đặc trưng tại mỗi nút phân chia. Thứ ba, quá trình huấn luyện các cây độc lập có thể thực hiện song song, giúp giảm đáng kể thời gian tính toán trên các hệ thống đa nhân.Nhờ các đặc điểm trên, Random Forest được ứng dụng rộng rãi trong y sinh học để phân loại dữ liệu gen và chẩn đoán bệnh, trong tài chính để phát hiện gian lận và đánh giá rủi ro tín dụng, cũng như trong các bài toán phân loại hình ảnh và xử lý ngôn ngữ tự nhiên.\\

\subsubsection*{Hạn chế} 

Bên cạnh các ưu điểm trên, Random Forest cũng tồn tại một số hạn chế cần lưu ý. Thứ nhất, do kết hợp hàng trăm cây quyết định, mô hình mất đi tính diễn giải trực quan vốn có của cây quyết định đơn lẻ, khó giải thích tại sao mô hình đưa ra một dự đoán cụ thể. Thứ hai, chi phí bộ nhớ và thời gian dự đoán tăng tuyến tính theo số lượng cây $B$, có thể trở thành hạn chế trong các hệ thống yêu cầu phản hồi thời gian thực. Cuối cùng, trên các tập dữ liệu có mất cân bằng lớp nghiêm trọng, mô hình có xu hướng thiên lệch về phía lớp đa số nếu không áp dụng các kỹ thuật bổ trợ phù hợp.

\subsection{XGBoost (Extreme Gradient Boosting)}

\subsection{MLP (Multilayer Perceptron)}
Perceptron đa tầng (Multilayer Perceptron - MLP) là một trong những kiến trúc mạng nơ-ron nhân tạo nền tảng và phổ biến nhất trong học máy hiện đại. Lấy cảm hứng từ cách thức hoạt động của hệ thần kinh sinh học, MLP mô phỏng quá trình truyền và xử lý thông tin qua các lớp tế bào nơ-ron nhân tạo được kết nối với nhau theo cấu trúc phân tầng. Mỗi đơn vị tính toán trong mạng tương tự một tế bào thần kinh, chúng nhận tín hiệu từ các đơn vị ở tầng trước, thực hiện một phép biến đổi có trọng số, sau đó truyền kết quả đến tầng tiếp theo.

Điểm then chốt phân biệt MLP so với các mô hình tuyến tính truyền thống nằm ở khả năng học các biểu diễn phi tuyến phức tạp trực tiếp từ dữ liệu. Trong khi các phương pháp như hồi quy logistic hay máy vector hỗ trợ tuyến tính (SVM) giả định rằng ranh giới quyết định có thể được xấp xỉ bằng một hàm tuyến tính của các đặc trưng đầu vào, MLP không đặt ra ràng buộc này. Thông qua việc xếp chồng nhiều tầng biến đổi phi tuyến, mô hình có thể học và biểu diễn các mẫu tương tác phức tạp giữa các đặc trưng mà không cần thiết kế thủ công bởi chuyên gia lĩnh vực. Tính chất này làm cho MLP trở thành một lựa chọn đặc biệt phù hợp cho các bài toán phân loại dữ liệu y tế, nơi mối quan hệ giữa các chỉ số lâm sàng và kết quả bệnh thường mang bản chất phi tuyến và phức tạp.

% ============================================================
\subsubsection{Kiến trúc cơ bản}
% ============================================================

Một mạng MLP điển hình được tổ chức thành ba loại tầng chức năng chính, mỗi tầng đóng vai trò xác định trong luồng xử lý thông tin của mô hình.

\textbf{Tầng đầu vào} (input layer) là tầng đầu tiên, có nhiệm vụ tiếp nhận dữ liệu thô từ bên ngoài. Số lượng đơn vị trong tầng này bằng đúng số chiều của vector đặc trưng đầu vào, tức là số lượng biến quan sát cho mỗi mẫu dữ liệu. Tầng đầu vào không thực hiện bất kỳ phép biến đổi nào mà đơn giản chuyển tiếp giá trị của từng đặc trưng đến tầng tiếp theo. Trong bối cảnh dữ liệu y tế, mỗi đơn vị đầu vào tương ứng với một chỉ số lâm sàng như huyết áp, nồng độ glucose, chỉ số khối cơ thể hay kết quả xét nghiệm máu.

\textbf{Các tầng ẩn} (hidden layers) là các tầng nằm giữa tầng đầu vào và tầng đầu ra, và là nơi phần lớn quá trình học biểu diễn diễn ra. Mỗi đơn vị trong tầng ẩn tính toán một tổ hợp có trọng số của tất cả đầu ra từ tầng trước đó, sau đó áp dụng một hàm phi tuyến lên kết quả. Qua nhiều tầng ẩn xếp chồng, mô hình dần dần xây dựng các biểu diễn trừu tượng hơn của dữ liệu đầu vào, từ các đặc trưng cơ bản ở tầng đầu đến các tổ hợp đặc trưng bậc cao phản ánh các mẫu phức tạp ở các tầng sâu hơn. Số lượng tầng ẩn và số đơn vị trong mỗi tầng là các siêu tham số kiến trúc ảnh hưởng trực tiếp đến khả năng biểu diễn và mức độ phức tạp của mô hình.

\textbf{Tầng đầu ra} (output layer) là tầng cuối cùng, chuyển đổi biểu diễn học được từ các tầng ẩn thành dạng phù hợp với bài toán cụ thể. Đối với bài toán phân loại nhị phân như chẩn đoán bệnh (có/không), tầng đầu ra thường có một đơn vị duy nhất với hàm kích hoạt sigmoid, trả về xác suất thuộc về lớp dương tính. Đối với phân loại đa lớp, tầng đầu ra có số đơn vị bằng số lớp, với hàm kích hoạt softmax đảm bảo tổng xác suất của tất cả các lớp bằng một.


\subsubsection{Học biểu diễn phi tuyến: Trọng số, bias và hàm kích hoạt}
% ============================================================

Khả năng học biểu diễn phi tuyến của MLP được tạo ra bởi sự kết hợp của ba thành phần cơ bản: trọng số liên kết, hệ số chệch và hàm kích hoạt. Mỗi kết nối giữa một đơn vị ở tầng trước và một đơn vị ở tầng sau được gắn với một giá trị trọng số $w$, biểu thị mức độ ảnh hưởng của tín hiệu đầu vào tương ứng. Ngoài ra, mỗi đơn vị trong mạng còn có một hệ số chệch $b$, hằng số cộng thêm trước khi áp dụng hàm kích hoạt cho phép mô hình dịch chuyển hàm kích hoạt, tăng tính linh hoạt trong việc khớp với dữ liệu. Khi được khởi tạo ngẫu nhiên và cập nhật trong quá trình huấn luyện, trọng số và hệ số chệch là các tham số mà mô hình học được từ dữ liệu.

Nếu chỉ dùng tổ hợp tuyến tính của các đầu vào, dù xếp chồng bao nhiêu tầng, mô hình cũng chỉ biểu diễn được các quan hệ tuyến tính vì tổ hợp tuyến tính của tổ hợp tuyến tính vẫn là tuyến tính. Hàm kích hoạt giải quyết vấn đề này bằng cách đưa vào tính phi tuyến sau mỗi lớp tổ hợp tuyến tính, cho phép mạng xấp xỉ bất kỳ hàm liên tục nào với độ chính xác tùy ý khi có đủ đơn vị ẩn --- tính chất được gọi là định lý xấp xỉ toàn năng (Cybenko, 1989; Hornik, 1991).

% ============================================================
\subsubsection{Nguyên lý hoạt động}
% ============================================================

Phép tính toán tại một đơn vị nơ-ron đơn lẻ diễn ra theo hai bước tuần tự. Đầu tiên, đơn vị tính tổ hợp tuyến tính của tất cả đầu vào mà nó nhận được:

\begin{equation}
  z = \sum_{i=1}^{n} w_i x_i + b
  = \mathbf{w}^T \mathbf{x} + b
  \label{eq:linear_combination}
\end{equation}

\noindent trong đó $x_i$ là giá trị của đầu vào thứ $i$ đến từ tầng trước (có thể là giá trị đặc trưng gốc nếu đây là tầng ẩn đầu tiên, hoặc đầu ra của đơn vị ở tầng ẩn trước), $w_i$ là trọng số tương ứng với kết nối từ đầu vào $x_i$ đến đơn vị hiện tại, $b$ là hệ số chệch của đơn vị này, $n$ là tổng số đầu vào mà đơn vị nhận được, và $z$ là giá trị tổng hợp chưa qua kích hoạt. Viết dưới dạng ma trận, $\mathbf{w}^T \mathbf{x}$ là tích vô hướng giữa vector trọng số $\mathbf{w} \in \mathbb{R}^n$ và vector đầu vào $\mathbf{x} \in \mathbb{R}^n$.

Bước thứ hai, giá trị $z$ được đưa qua hàm kích hoạt $\varphi(\cdot)$ để tạo ra đầu ra của đơn vị:

\begin{equation}
  a = \varphi(z) = \varphi\!\left(\mathbf{w}^T \mathbf{x} + b\right)
  \label{eq:activation}
\end{equation}

\noindent trong đó $a$ là đầu ra của đơn vị nơ-ron, sẽ được truyền đến tất cả các đơn vị ở tầng tiếp theo. Tổng quát hóa lên toàn bộ một tầng ẩn gồm $m$ đơn vị, lan truyền tiến từ tầng $l-1$ sang tầng $l$ được biểu diễn dưới dạng ma trận:

\begin{equation}
  \mathbf{a}^{(l)}
  = \varphi\!\left(\mathbf{W}^{(l)}\, \mathbf{a}^{(l-1)} + \mathbf{b}^{(l)}\right)
  \label{eq:layer_forward}
\end{equation}

\noindent trong đó $\mathbf{W}^{(l)} \in \mathbb{R}^{m_l \times m_{l-1}}$ là ma trận trọng số của tầng $l$ với $m_l$ là số đơn vị ở tầng $l$ và $m_{l-1}$ là số đơn vị ở tầng $l-1$; $\mathbf{b}^{(l)} \in \mathbb{R}^{m_l}$ là vector hệ số chệch; $\mathbf{a}^{(l-1)}$ là vector đầu ra từ tầng trước; và $\varphi(\cdot)$ được áp dụng theo từng phần tử của vector. Phép tính này được lặp lại từ tầng đầu vào qua tất cả các tầng ẩn cho đến tầng đầu ra.

\subsubsection{Vai trò của hàm kích hoạt}

Hàm kích hoạt là thành phần quyết định tính phi tuyến của mạng. Nếu không có hàm kích hoạt, dù xếp chồng bao nhiêu tầng, toàn bộ mạng vẫn chỉ tương đương một phép biến đổi tuyến tính. Ba hàm kích hoạt được sử dụng phổ biến nhất trong kiến trúc MLP hiện đại là ReLU, sigmoid và softmax, mỗi hàm phù hợp với một vị trí và mục đích cụ thể trong mạng.\\

\textbf{Hàm đơn vị tuyến tính có giới hạn dưới} (Rectified Linear Unit - ReLU) được định nghĩa là:
\begin{equation}
    \varphi_{\text{ReLU}}(z) = \max(0, z)
    \label{eq:relu}
\end{equation}
\noindent trong đó $z$ là giá trị tổng hợp tuyến tính nhận được từ tầng trước, tức là đầu vào chưa qua kích hoạt của một đơn vị nơ-ron. Hàm này hoạt động theo nguyên tắc đơn giản: giữ nguyên giá trị nếu $z > 0$, và trả về không nếu $z \leq 0$. Về mặt đạo hàm, $\varphi'_{\text{ReLU}}(z) = 1$ khi $z > 0$ và bằng không khi $z < 0$, nghĩa là gradient được truyền nguyên vẹn qua các nơ-ron đang hoạt động thay vì bị nén về không như ở các hàm kích hoạt cổ điển. Tính chất này giúp tránh hiện tượng suy giảm gradient, làm cho quá trình huấn luyện mạng sâu trở nên ổn định và hiệu quả hơn. Nhờ đó, ReLU là lựa chọn mặc định cho các tầng ẩn trong hầu hết các kiến trúc mạng nơ-ron hiện đại.

\textbf{Hàm softmax} là tổng quát hóa của hàm sigmoid cho bài toán phân loại đa lớp. Với vector đầu vào $\mathbf{z} = (z_1, z_2, \ldots, z_K)$ là đầu ra tuyến tính của tầng cuối cùng, đầu ra softmax tại vị trí $k$ được tính là:
\begin{equation}
    \varphi_{\text{softmax}}(z_k) = \frac{e^{z_k}}{\displaystyle\sum_{j=1}^{K} e^{z_j}}, \quad k = 1, \ldots, K
    \label{eq:softmax}
\end{equation}
trong đó $K$ là tổng số lớp, $z_k$ là giá trị logit tương ứng với lớp $k$, và mẫu số $\sum_{j=1}^{K} e^{z_j}$ là tổng của tất cả các giá trị mũ nhằm mục đích chuẩn hóa. Phép chuẩn hóa này đảm bảo hai tính chất quan trọng: mọi đầu ra đều dương, và tổng tất cả đầu ra bằng một, tức là tạo thành một phân phối xác suất hợp lệ trên $K$ lớp. Lớp có giá trị $\varphi_{\text{softmax}}(z_k)$ lớn nhất được chọn làm nhãn dự đoán cuối cùng. Trong khóa luận này, softmax được sử dụng tại tầng đầu ra của MLP cho cả hai bài toán phân loại đa lớp.\\

% ============================================================
\subsubsection{Vai trò của hàm kích hoạt}
% ============================================================

Hàm kích hoạt là thành phần quyết định tính phi tuyến của mạng và ảnh hưởng trực tiếp đến hiệu quả của quá trình huấn luyện. Trong lịch sử phát triển của mạng nơ-ron, nhiều hàm kích hoạt khác nhau đã được đề xuất và sử dụng, mỗi hàm có đặc tính phù hợp với các ngữ cảnh khác nhau trong kiến trúc mạng.

\textbf{Hàm đơn vị tuyến tính có giới chặn dưới} (\textit{Rectified Linear Unit} --- ReLU) là hàm kích hoạt phổ biến nhất cho các tầng ẩn trong các mạng nơ-ron hiện đại, được định nghĩa là:

\begin{equation}
  \varphi_{\mathrm{ReLU}}(z) = \max(0, z)
  \label{eq:relu}
\end{equation}

\noindent Hàm này giữ nguyên giá trị dương và đưa mọi giá trị âm về không. Sự đơn giản về mặt tính toán của ReLU --- đạo hàm bằng một khi $z > 0$ và bằng không khi $z < 0$ --- giúp tránh được hiện tượng suy giảm gradient (\textit{vanishing gradient}) thường gặp với các hàm kích hoạt cổ điển hơn, làm cho quá trình huấn luyện mạng sâu trở nên ổn định và hiệu quả hơn đáng kể.

\textbf{Hàm sigmoid} ép mọi giá trị thực vào khoảng $(0, 1)$:

\begin{equation}
  \varphi_{\mathrm{sigmoid}}(z) = \frac{1}{1 + e^{-z}}
  \label{eq:sigmoid}
\end{equation}

\noindent Nhờ miền giá trị nằm trong $(0, 1)$, đầu ra của hàm sigmoid có thể được diễn giải tự nhiên như xác suất. Do đó, hàm sigmoid thường được sử dụng tại đơn vị đầu ra duy nhất trong bài toán phân loại nhị phân --- chẳng hạn, xác suất một bệnh nhân mắc bệnh tiểu đường hay không.

\textbf{Hàm softmax} là tổng quát hóa của hàm sigmoid cho bài toán phân loại đa lớp. Với vector đầu vào $\mathbf{z} = (z_1, z_2, \ldots, z_K)$, đầu ra của softmax tại vị trí $k$ là:

\begin{equation}
  \varphi_{\mathrm{softmax}}(z_k)
  = \frac{e^{z_k}}{\displaystyle\sum_{j=1}^{K} e^{z_j}}, \quad k = 1, \ldots, K
  \label{eq:softmax}
\end{equation}

\noindent trong đó $K$ là tổng số lớp, và đầu ra $\varphi_{\mathrm{softmax}}(z_k)$ là xác suất ước lượng của mẫu thuộc lớp $k$. Hàm softmax đảm bảo hai tính chất quan trọng: mọi đầu ra đều dương, và tổng tất cả đầu ra bằng một, tức là tạo thành một phân phối xác suất hợp lệ trên $K$ lớp.

% ============================================================
\subsubsection{Quá trình huấn luyện}

Huấn luyện MLP là quá trình tìm bộ trọng số và hệ số chệch sao cho mô hình đưa ra dự đoán gần nhất với nhãn thực tế trên tập dữ liệu huấn luyện. Quá trình này diễn ra qua ba giai đoạn liên tiếp, được lặp lại nhiều lần cho đến khi mô hình hội tụ: lan truyền tiến, tính hàm mất mát, và lan truyền ngược.\\

\textbf{Lan truyền tiến}

Lan truyền tiến là quá trình tính toán đầu ra của mạng cho một mẫu đầu vào $\mathbf{x}$. Bắt đầu từ tầng đầu vào, tín hiệu được truyền lần lượt qua từng tầng theo công thức:

\begin{equation}
    \mathbf{a}^{(l)} = \varphi\!\left(\mathbf{W}^{(l)}\mathbf{a}^{(l-1)} + \mathbf{b}^{(l)}\right)
    \label{eq:forward}
\end{equation}

\noindent trong đó $\mathbf{W}^{(l)} \in \mathbb{R}^{m_l \times m_{l-1}}$ là ma trận trọng số của tầng $l$, với $m_l$ và $m_{l-1}$ lần lượt là số đơn vị ở tầng $l$ và tầng $l-1$; $\mathbf{b}^{(l)} \in \mathbb{R}^{m_l}$ là vector hệ số chệch; $\mathbf{a}^{(l-1)}$ là vector đầu ra từ tầng trước, với $\mathbf{a}^{(0)} = \mathbf{x}$ là vector đặc trưng đầu vào; và $\varphi(\cdot)$ là hàm kích hoạt được áp dụng theo từng phần tử. Phép tính này được lặp lại từ tầng đầu vào qua tất cả các tầng ẩn cho đến tầng đầu ra, nơi mô hình tạo ra dự đoán $\hat{y}$.

\textbf{Hàm mất mát}

Sau khi có dự đoán $\hat{y}$, sai lệch giữa dự đoán và nhãn thực tế $y$ được đo bằng hàm mất mát. Đối với bài toán phân loại nhị phân, hàm mất mát thường dùng là entropy chéo nhị phân, được định nghĩa là:
\begin{equation}
    \mathcal{L}(y, \hat{y}) = -\left[y \log \hat{y} + (1 - y) \log(1 - \hat{y})\right]
    \label{eq:bce}
\end{equation}
trong đó $y \in \{0, 1\}$ là nhãn thực tế và $\hat{y} \in (0, 1)$ là xác suất dự đoán mẫu thuộc lớp dương tính. Hàm mất mát này phạt nặng các dự đoán tự tin nhưng sai: khi $y = 1$ mà $\hat{y}$ gần $0$, mất mát tiến đến vô cùng; ngược lại, mất mát tiến về không khi dự đoán đúng với độ tự tin cao. Trên toàn bộ tập dữ liệu gồm $n$ mẫu, mất mát tổng được lấy trung bình:
\begin{equation}
    \mathcal{L}_{\text{total}} = \frac{1}{n} \sum_{i=1}^{n} \mathcal{L}(y_i, \hat{y}_i)
    \label{eq:total_loss}
\end{equation}

\textbf{Lan truyền ngược}

Để cập nhật tham số nhằm giảm mất mát, mô hình cần tính đạo hàm của hàm mất mát theo từng tham số. Lan truyền ngược là thuật toán tính hiệu quả các đạo hàm này bằng cách áp dụng quy tắc dây chuyền của vi phân, truyền ngược từ tầng đầu ra về tầng đầu vào. Tại tầng $l$, gradient của hàm mất mát theo ma trận trọng số được tính là:
\begin{equation}
    \frac{\partial \mathcal{L}}{\partial \mathbf{W}^{(l)}} = \boldsymbol{\delta}^{(l)} \cdot \left(\mathbf{a}^{(l-1)}\right)^T
    \label{eq:grad_W}
\end{equation}
trong đó $\boldsymbol{\delta}^{(l)}$ là vector tín hiệu lỗi ngược tại tầng $l$, được tính đệ quy từ gradient tại tầng $l+1$:
\begin{equation}
    \boldsymbol{\delta}^{(l)} = \left(\mathbf{W}^{(l+1)}\right)^T \boldsymbol{\delta}^{(l+1)} \odot \varphi'\!\left(\mathbf{z}^{(l)}\right)
    \label{eq:delta}
\end{equation}
trong đó $\odot$ là phép nhân theo từng phần tử (Hadamard), $\mathbf{z}^{(l)} = \mathbf{W}^{(l)}\mathbf{a}^{(l-1)} + \mathbf{b}^{(l)}$ là giá trị tổng hợp tuyến tính trước kích hoạt tại tầng $l$, và $\varphi'(\cdot)$ là đạo hàm của hàm kích hoạt. Chuỗi tính toán đệ quy này cho phép truyền ngược lỗi từ tầng đầu ra qua tất cả các tầng ẩn chỉ trong một lần duyệt, đảm bảo hiệu quả tính toán ngay cả với mạng rất sâu.

\textbf{Tối ưu hóa bằng hạ gradient}

Sau khi có gradient, các tham số được cập nhật theo hướng giảm mất mát. Quy tắc cập nhật cơ bản theo phương pháp hạ gradient ngẫu nhiên là:
\begin{equation}
    \mathbf{W}^{(l)} \leftarrow \mathbf{W}^{(l)} - \eta \cdot \frac{\partial \mathcal{L}}{\partial \mathbf{W}^{(l)}}
    \label{eq:sgd}
\end{equation}
trong đó $\eta > 0$ là tốc độ học, kiểm soát kích thước bước di chuyển trong không gian tham số mỗi lần cập nhật. Trong thực tế, thuật toán Adam thường được ưa dùng hơn SGD thuần túy, do Adam kết hợp động lượng và tốc độ học thích nghi cho từng tham số, giúp hội tụ nhanh hơn và ổn định hơn trên nhiều loại bài toán. Toàn bộ chu kỳ lan truyền tiến $\rightarrow$ tính mất mát $\rightarrow$ lan truyền ngược $\rightarrow$ cập nhật tham số được lặp lại nhiều lần trên toàn bộ tập huấn luyện cho đến khi mất mát hội tụ về một giá trị đủ nhỏ hoặc hiệu năng trên tập kiểm định không còn cải thiện.

\subsubsection{Ưu điểm của MLP trong bài toán dữ kiệu y tế}
% ============================================================

Đối với bài toán phân loại dữ liệu chăm sóc sức khỏe, MLP mang lại một số ưu điểm quan trọng so với các mô hình truyền thống. Trước hết, dữ liệu y tế thường chứa đựng các mối quan hệ phi tuyến phức tạp giữa các chỉ số lâm sàng và kết quả bệnh: ví dụ, tác động của nồng độ glucose máu đến nguy cơ mắc tiểu đường không phải là tuyến tính đơn giản mà có sự tương tác với các yếu tố khác như chỉ số khối cơ thể và tuổi tác. MLP có khả năng tự động phát hiện và mô hình hóa các mối quan hệ và tương tác đặc trưng bậc cao này mà không cần chuyên gia phải thiết kế thủ công, nhờ vào chuỗi biến đổi phi tuyến qua các tầng ẩn.\\

Thứ hai, MLP không đặt ra các giả định phân phối về dữ liệu đầu vào, không yêu cầu phân phối chuẩn, không cần tính độc lập giữa các đặc trưng, điều này đặc biệt có giá trị trong dữ liệu y tế thực tế vốn thường không thỏa mãn các giả định thống kê lý tưởng. Thứ ba, MLP có khả năng mở rộng tự nhiên sang bài toán phân loại đa lớp thông qua tầng đầu ra với hàm softmax, phù hợp cho các bài toán phân loại nhiều mức độ bệnh. Cuối cùng, với kích thước mạng phù hợp, MLP có khả năng biểu diễn các hàm quyết định rất phức tạp, tiệm cận khả năng của các kiến trúc học sâu hiện đại trong khi vẫn duy trì được sự đơn giản về mặt huấn luyện và triển khai.\\

% ============================================================
\subsubsection{Hạn chế}
% ============================================================

Mặc dù có nhiều ưu điểm, MLP cũng đặt ra một số thách thức cần được xem xét cẩn thận trong thực tế.\\

Xu hướng quá khớp dữ liệu huấn luyện: Với số lượng tham số lớn, một mạng MLP vừa có thể có hàng nghìn đến hàng triệu tham số khiến mô hình dễ ``ghi nhớ'' chi tiết và nhiễu của tập huấn luyện thay vì học các mẫu tổng quát, dẫn đến hiệu năng kém trên dữ liệu mới. Các kỹ thuật như \textit{dropout} (tắt ngẫu nhiên một tỉ lệ nơ-ron trong mỗi bước huấn luyện), chính quy hóa L2 trên trọng số và dừng sớm khi hiệu năng kiểm định không còn cải thiện là các biện pháp kiểm soát quá khớp được sử dụng rộng rãi. Thách thức này càng trở nên nổi bật hơn khi tập dữ liệu y tế thường có kích thước hạn chế so với số lượng đặc trưng.\\

Phụ thuộc vào chuẩn hóa dữ liệu: khác với cây quyết định hay rừng ngẫu nhiên, các phương pháp chỉ sử dụng thứ tự và ngưỡng của giá trị đặc trưng, MLP thực hiện các phép nhân và cộng trực tiếp trên giá trị số. Điều này làm cho mô hình rất nhạy cảm với phạm vi giá trị của các đặc trưng: nếu một đặc trưng có giá trị trong khoảng [0, 1000] trong khi đặc trưng khác có giá trị trong [0, 1], các trọng số liên kết với đặc trưng thứ nhất sẽ có xu hướng chi phối quá mức quá trình học, làm lệch gradient và cản trở hội tụ. Do đó, chuẩn hóa dữ liệu đầu vào, thường là chuẩn hóa về phạm vi $[0, 1]$ hoặc về phân phối có trung bình bằng không và độ lệch chuẩn bằng một là bước tiền xử lý bắt buộc trước khi huấn luyện MLP.\\

Số lượng siêu tham số lớn và nhạy cảm: hiệu năng của MLP phụ thuộc vào nhiều quyết định thiết kế: số tầng ẩn, số đơn vị trong mỗi tầng, loại hàm kích hoạt, tốc độ học, kích thước batch, tỉ lệ dropout, và các hệ số chính quy hóa. Không có một bộ siêu tham số tối ưu nào phù hợp với mọi bài toán, và quá trình tìm kiếm siêu tham số tốt, thường thông qua tìm kiếm lưới hoặc tìm kiếm ngẫu nhiên kết hợp kiểm định chéo, đòi hỏi thời gian và tài nguyên tính toán đáng kể.\\

Khó diễn giải kết quả: không giống như hồi quy logistic, nơi hệ số của từng biến trực tiếp phản ánh ảnh hưởng của biến đó lên kết quả hay cây quyết định có thể được đọc như tập quy tắc rõ ràng, MLP là một ``hộp đen'' về mặt diễn giải. Với hàng nghìn tham số phân tán qua nhiều tầng, việc trả lời câu hỏi "tại sao mô hình đưa ra dự đoán này cho bệnh nhân cụ thể đó?" đòi hỏi các kỹ thuật diễn giải bổ sung như SHAP values hay LIME. Hạn chế này có thể là rào cản đáng kể trong các ứng dụng y tế, nơi tính minh bạch và khả năng giải thích của quyết định lâm sàng là yêu cầu quan trọng.\\

% ============================================================